\documentclass[11pt]{article}
% The preamble every project's paper shares: % The preamble every project's paper shares: % The preamble every project's paper shares: \input{../../preamble} from
% projects/<project>/paper/. Boards come from chessboard and xskak, plots from
% pgfplots; every number a paper prints comes from its generated.tex.
\usepackage[LSBC4,T1]{fontenc}
\usepackage[utf8]{inputenc}
\usepackage{lmodern}
\usepackage[margin=1in]{geometry}
\usepackage{microtype}
\usepackage{amsmath}
\usepackage{booktabs}
\usepackage{tabularx}
\usepackage{array}
\usepackage{xcolor}
\usepackage[obeyspaces,spaces]{url}
\usepackage{xskak}
\usepackage{pgfplots}
\pgfplotsset{compat=1.18}
\usepackage{caption}
\usepackage{subcaption}
\usepackage{enumitem}
\usepackage[section]{placeins}
\usepackage[hidelinks]{hyperref}

\setchessboard{boardfontencoding=LSBC4, boardfontfamily=skaknew, blackfieldcolor=black!22,
  whitefieldcolor=white, setfontcolors, showmover=false, label=true, labelfontsize=6pt,
  boardfontsize=14pt}

% A move or line in UCI, and a position in FEN, both breakable at a slash.
\newcommand{\uci}[1]{\texttt{#1}}
\newcommand{\fen}[1]{{\footnotesize\path{#1}}}
\newcolumntype{L}{>{\raggedright\arraybackslash}X}

\setlength{\parskip}{0.4em}
\setlength{\parindent}{0pt}
 from
% projects/<project>/paper/. Boards come from chessboard and xskak, plots from
% pgfplots; every number a paper prints comes from its generated.tex.
\usepackage[LSBC4,T1]{fontenc}
\usepackage[utf8]{inputenc}
\usepackage{lmodern}
\usepackage[margin=1in]{geometry}
\usepackage{microtype}
\usepackage{amsmath}
\usepackage{booktabs}
\usepackage{tabularx}
\usepackage{array}
\usepackage{xcolor}
\usepackage[obeyspaces,spaces]{url}
\usepackage{xskak}
\usepackage{pgfplots}
\pgfplotsset{compat=1.18}
\usepackage{caption}
\usepackage{subcaption}
\usepackage{enumitem}
\usepackage[section]{placeins}
\usepackage[hidelinks]{hyperref}

\setchessboard{boardfontencoding=LSBC4, boardfontfamily=skaknew, blackfieldcolor=black!22,
  whitefieldcolor=white, setfontcolors, showmover=false, label=true, labelfontsize=6pt,
  boardfontsize=14pt}

% A move or line in UCI, and a position in FEN, both breakable at a slash.
\newcommand{\uci}[1]{\texttt{#1}}
\newcommand{\fen}[1]{{\footnotesize\path{#1}}}
\newcolumntype{L}{>{\raggedright\arraybackslash}X}

\setlength{\parskip}{0.4em}
\setlength{\parindent}{0pt}
 from
% projects/<project>/paper/. Boards come from chessboard and xskak, plots from
% pgfplots; every number a paper prints comes from its generated.tex.
\usepackage[LSBC4,T1]{fontenc}
\usepackage[utf8]{inputenc}
\usepackage{lmodern}
\usepackage[margin=1in]{geometry}
\usepackage{microtype}
\usepackage{amsmath}
\usepackage{booktabs}
\usepackage{tabularx}
\usepackage{array}
\usepackage{xcolor}
\usepackage[obeyspaces,spaces]{url}
\usepackage{xskak}
\usepackage{pgfplots}
\pgfplotsset{compat=1.18}
\usepackage{caption}
\usepackage{subcaption}
\usepackage{enumitem}
\usepackage[section]{placeins}
\usepackage[hidelinks]{hyperref}

\setchessboard{boardfontencoding=LSBC4, boardfontfamily=skaknew, blackfieldcolor=black!22,
  whitefieldcolor=white, setfontcolors, showmover=false, label=true, labelfontsize=6pt,
  boardfontsize=14pt}

% A move or line in UCI, and a position in FEN, both breakable at a slash.
\newcommand{\uci}[1]{\texttt{#1}}
\newcommand{\fen}[1]{{\footnotesize\path{#1}}}
\newcolumntype{L}{>{\raggedright\arraybackslash}X}

\setlength{\parskip}{0.4em}
\setlength{\parindent}{0pt}

% Written by projects/judge-factors/paper/generate.mjs; do not edit by hand.

\newcommand{\AppCommit}{f4bd63f19b}
\newcommand{\LeanCp}{80}
\newcommand{\TwiceLeanCp}{160}
\newcommand{\WorstLean}{three}
\newcommand{\TwiceWorstLean}{six}
\newcommand{\LevelCp}{50}
\newcommand{\ClearCp}{150}
\newcommand{\PlyFloor}{12}
\newcommand{\TacticBar}{one}
\newcommand{\BudgetNodes}{250,000}
\newcommand{\Cap}{3,000}
\newcommand{\MinHeldout}{200}
\newcommand{\ReliablePct}{67}
\newcommand{\BandLow}{1250--1449}
\newcommand{\BandHigh}{1450--1649}
\newcommand{\RatingLow}{1250}
\newcommand{\RatingHigh}{1649}
\newcommand{\WShield}{25}
\newcommand{\WRingAttacker}{12}
\newcommand{\WHalfOpenFile}{20}
\newcommand{\WOpenFile}{30}
\newcommand{\WMobility}{4}
\newcommand{\WSpace}{4}
\newcommand{\WDoubled}{15}
\newcommand{\WIsolated}{20}
\newcommand{\WPassed}{20}
\newcommand{\WPassedRank}{12}
\newcommand{\GamesPerMonth}{8,000}
\newcommand{\Minutes}{10}
\newcommand{\MinPlies}{30}
\newcommand{\TrainRead}{247,185}
\newcommand{\HeldRead}{215,414}
\newcommand{\TrainRapid}{34,809}
\newcommand{\HeldRapid}{35,069}
\newcommand{\TrainPositions}{566,133}
\newcommand{\HeldPositions}{570,663}
\newcommand{\TrainLater}{470,133}
\newcommand{\HeldLater}{474,663}
\newcommand{\TrainQualifying}{103,668}
\newcommand{\HeldQualifying}{105,341}
\newcommand{\TrainQualifyingPct}{22}
\newcommand{\HeldQualifyingPct}{22}
\newcommand{\TrainLevelDrop}{62}
\newcommand{\HeldLevelDrop}{62}
\newcommand{\TrainTactical}{28,750}
\newcommand{\HeldTactical}{29,422}
\newcommand{\Rows}{23,253}
\newcommand{\Unique}{21,266}
\newcommand{\KingTrainKept}{1,852}
\newcommand{\KingHeldKept}{1,919}
\newcommand{\KingTrainLeans}{1,852}
\newcommand{\KingHeldLeans}{1,919}
\newcommand{\KingVerdict}{tells who is better}
\newcommand{\KingTrain}{77.2 (74.1--80.1)}
\newcommand{\KingTrainN}{747}
\newcommand{\KingTrainPct}{77.2}
\newcommand{\KingTrainLevel}{29}
\newcommand{\KingTrainClear}{40}
\newcommand{\KingHeld}{81.2 (78.1--83.8)}
\newcommand{\KingHeldN}{732}
\newcommand{\KingHeldPct}{81.2}
\newcommand{\KingHeldLevel}{30}
\newcommand{\KingHeldClear}{38}
\newcommand{\ActivityTrainKept}{3,000}
\newcommand{\ActivityHeldKept}{3,000}
\newcommand{\ActivityTrainLeans}{3,667}
\newcommand{\ActivityHeldLeans}{3,634}
\newcommand{\ActivityVerdict}{tells who is better}
\newcommand{\ActivityTrain}{81.7 (78.8--84.3)}
\newcommand{\ActivityTrainN}{776}
\newcommand{\ActivityTrainPct}{81.7}
\newcommand{\ActivityTrainLevel}{35}
\newcommand{\ActivityTrainClear}{26}
\newcommand{\ActivityHeld}{82.3 (79.6--84.8)}
\newcommand{\ActivityHeldN}{826}
\newcommand{\ActivityHeldPct}{82.3}
\newcommand{\ActivityHeldLevel}{37}
\newcommand{\ActivityHeldClear}{28}
\newcommand{\StructureTrainKept}{723}
\newcommand{\StructureHeldKept}{759}
\newcommand{\StructureTrainLeans}{723}
\newcommand{\StructureHeldLeans}{759}
\newcommand{\StructureVerdict}{tells who is better}
\newcommand{\StructureTrain}{63.0 (58.1--67.7)}
\newcommand{\StructureTrainN}{384}
\newcommand{\StructureTrainPct}{63.0}
\newcommand{\StructureTrainLevel}{26}
\newcommand{\StructureTrainClear}{53}
\newcommand{\StructureHeld}{68.9 (64.2--73.2)}
\newcommand{\StructureHeldN}{405}
\newcommand{\StructureHeldPct}{68.9}
\newcommand{\StructureHeldLevel}{25}
\newcommand{\StructureHeldClear}{53}
\newcommand{\LooseTrainKept}{3,000}
\newcommand{\LooseHeldKept}{3,000}
\newcommand{\LooseTrainLeans}{4,548}
\newcommand{\LooseHeldLeans}{4,599}
\newcommand{\LooseVerdict}{tells who is better}
\newcommand{\LooseTrain}{92.6 (91.1--93.9)}
\newcommand{\LooseTrainN}{1,409}
\newcommand{\LooseTrainPct}{92.6}
\newcommand{\LooseTrainLevel}{21}
\newcommand{\LooseTrainClear}{47}
\newcommand{\LooseHeld}{92.9 (91.5--94.2)}
\newcommand{\LooseHeldN}{1,388}
\newcommand{\LooseHeldPct}{92.9}
\newcommand{\LooseHeldLevel}{20}
\newcommand{\LooseHeldClear}{46}
\newcommand{\WorstTrainKept}{3,000}
\newcommand{\WorstHeldKept}{3,000}
\newcommand{\WorstTrainLeans}{3,168}
\newcommand{\WorstHeldLeans}{3,206}
\newcommand{\WorstVerdict}{tells who is better}
\newcommand{\WorstTrain}{71.5 (68.7--74.1)}
\newcommand{\WorstTrainN}{1,055}
\newcommand{\WorstTrainPct}{71.5}
\newcommand{\WorstTrainLevel}{34}
\newcommand{\WorstTrainClear}{35}
\newcommand{\WorstHeld}{74.2 (71.4--76.8)}
\newcommand{\WorstHeldN}{1,023}
\newcommand{\WorstHeldPct}{74.2}
\newcommand{\WorstHeldLevel}{35}
\newcommand{\WorstHeldClear}{34}
\newcommand{\RuleTable}{%
\begin{table}[ht]\centering\small
\setlength{\tabcolsep}{3pt}
\begin{tabular}{llrrrrll}
\toprule
 & & \multicolumn{2}{c}{Train} & \multicolumn{2}{c}{Held out} & & \\
\cmidrule(lr){3-4}\cmidrule(lr){5-6}
Reading & Leans at & Clear & Agreement, \% & Clear & Agreement, \% & Verdict & Reliable \\
\midrule
King safety & 80 cp & 747 & 77.2 (74.1--80.1) & 732 & 81.2 (78.1--83.8) & tells who is better & yes \\
Activity and space & 80 cp & 776 & 81.7 (78.8--84.3) & 826 & 82.3 (79.6--84.8) & tells who is better & yes \\
Pawn structure & 80 cp & 384 & 63.0 (58.1--67.7) & 405 & 68.9 (64.2--73.2) & tells who is better & no \\
Something loose & 80 cp & 1,409 & 92.6 (91.1--93.9) & 1,388 & 92.9 (91.5--94.2) & tells who is better & yes \\
The worst piece & 3 moves & 1,055 & 71.5 (68.7--74.1) & 1,023 & 74.2 (71.4--76.8) & tells who is better & yes \\
\bottomrule
\end{tabular}
\caption{Each reading's agreement with the engine: of its positions with a clear verdict (\emph{Clear}), the share where the side the reading favours is the side Stockfish says is better, with the 95\% Wilson interval. The verdict, and whether the reading is reliable on its own, are the rule's, fixed before any position was scored.}
\label{tab:rule}
\end{table}
}
\newcommand{\DescTable}{%
\begin{table}[ht]\centering\small
\setlength{\tabcolsep}{3pt}
\begin{tabular}{lrrrrrrrr}
\toprule
 & \multicolumn{2}{c}{Engine's verdict} & \multicolumn{2}{c}{Band} & \multicolumn{3}{c}{Phase} & \\
\cmidrule(lr){2-3}\cmidrule(lr){4-5}\cmidrule(lr){6-8}
Reading & Level & Clear & 1250--1449 & 1450--1649 & Opening & Middle & End & Twice the lean \\
\midrule
King safety & 30 & 38 & 81 (340) & 82 (392) & 78 (136) & 82 (541) & 82 (55) & 97 (32) \\
Activity and space & 37 & 28 & 80 (389) & 84 (437) & 79 (248) & 84 (473) & 84 (105) & -- (18) \\
Pawn structure & 25 & 53 & 65 (182) & 72 (223) & -- (6) & 65 (122) & 71 (277) & 69 (49) \\
Something loose & 20 & 46 & 94 (621) & 92 (767) & 92 (338) & 93 (737) & 94 (313) & -- (0) \\
The worst piece & 35 & 34 & 75 (449) & 74 (574) & -- (11) & 76 (630) & 71 (382) & 74 (148) \\
\bottomrule
\end{tabular}
\caption{Held-out positions only. \emph{Level} and \emph{Clear}: the share of a reading's positions the engine calls level or clear. The other columns give the agreement, with the number of clear positions in brackets, by the band of the two players' mean rating, by the divider's phase, and where the reading is at least twice its lean; a cell with fewer than 30 clear positions shows only its count. These were named before the run and carry no rule.}
\label{tab:desc}
\end{table}
}
\newcommand{\ResultTable}{%
\begin{table}[ht]\centering\small
\begin{tabular}{lrrr}
\toprule
Reading & Decisive games & Reading's side won & Engine's side won \\
\midrule
King safety & 700 & 64 (60--67) & 69 (65--72) \\
Activity and space & 790 & 64 (61--68) & 70 (67--73) \\
Pawn structure & 375 & 65 (60--69) & 78 (74--82) \\
Something loose & 1,318 & 67 (64--69) & 70 (67--72) \\
The worst piece & 975 & 63 (60--66) & 72 (69--75) \\
\bottomrule
\end{tabular}
\caption{Held-out positions with a clear verdict from games that ended in a win: how often the side the reading favours won the game, beside how often the side the engine favours won it, with 95\% Wilson intervals.}
\label{tab:result}
\end{table}
}
\newcommand{\SumTrain}{82.5 (81.3--83.7)}
\newcommand{\SumHeld}{83.8 (82.6--84.9)}
\newcommand{\SumTrainN}{3,838}
\newcommand{\SumHeldN}{3,813}
\newcommand{\LogTrain}{84.5}
\newcommand{\LogHeld}{85.6}
\newcommand{\LogTrainN}{3,848}
\newcommand{\LogHeldN}{3,820}
\newcommand{\KingWeight}{0.96}
\newcommand{\ActivityWeight}{1.48}
\newcommand{\StructureWeight}{0.52}
\newcommand{\LooseWeight}{2.41}
\newcommand{\WorstWeight}{0.06}
\newcommand{\Ranked}{something loose (\LooseHeldPct\%), activity and space (\ActivityHeldPct\%), king safety (\KingHeldPct\%), the worst piece (\WorstHeldPct\%), then pawn structure (\StructureHeldPct\%)}
\newcommand{\LevelMin}{20}
\newcommand{\LevelMax}{37}
\newcommand{\KingTwiceHeld}{97}
\newcommand{\KingTwiceHeldN}{32}
\newcommand{\ActivityTwiceHeld}{94}
\newcommand{\ActivityTwiceHeldN}{18}
\newcommand{\StructureTwiceHeld}{69}
\newcommand{\StructureTwiceHeldN}{49}
\newcommand{\WorstTwiceHeld}{74}
\newcommand{\WorstTwiceHeldN}{148}
\newcommand{\WonMin}{63}
\newcommand{\WonMax}{67}
\newcommand{\StructureWon}{65}
\newcommand{\StructureEngineWon}{78}
\newcommand{\MoverTrain}{64.1 (62.0--66.1)}
\newcommand{\MoverHeld}{65.8 (63.7--67.8)}
\newcommand{\MoverHeldN}{2,074}
\newcommand{\WhiteHeld}{58}
\newcommand{\MoverHeldPct}{66}
\newcommand{\ForMoverMin}{90}
\newcommand{\ForMoverMax}{98}
\newcommand{\KingForOther}{61}
\newcommand{\KingForOtherTrain}{58}
\newcommand{\KingAlone}{78}
\newcommand{\ActivityForOther}{65}
\newcommand{\ActivityForOtherTrain}{65}
\newcommand{\ActivityAlone}{82}
\newcommand{\StructureForOther}{45}
\newcommand{\StructureForOtherTrain}{39}
\newcommand{\StructureAlone}{75}
\newcommand{\LooseForOther}{--}
\newcommand{\LooseForOtherTrain}{--}
\newcommand{\LooseAlone}{95}
\newcommand{\WorstForOther}{51}
\newcommand{\WorstForOtherTrain}{51}
\newcommand{\WorstAlone}{74}
\newcommand{\SideTable}{%
\begin{table}[ht]\centering\small
\begin{tabular}{lrrrr}
\toprule
 & & \multicolumn{3}{c}{Agreement, \%} \\
\cmidrule(lr){3-5}
Reading & Favours the mover & For the mover & For the other side & Only sign \\
\midrule
King safety & 42 & 96 (93--97) & 61 (55--66) & 78 (74--82) \\
Activity and space & 38 & 98 (96--99) & 65 (60--69) & 82 (79--85) \\
Pawn structure & 47 & 90 (85--93) & 45 (38--52) & 75 (69--80) \\
Something loose & 100 & 93 (91--94) & -- & 95 (94--97) \\
The worst piece & 43 & 93 (91--95) & 51 (47--56) & 74 (69--78) \\
\bottomrule
\end{tabular}
\caption{Added after the run; held-out positions. \emph{Favours the mover}: the share of a reading's positions where it favours the side to move. The agreement, with 95\% Wilson intervals, when it favours the side to move, when it favours the side that just moved, and when no other drill factor leans.}
\label{tab:side}
\end{table}
}
\newcommand{\LooseOverStructure}{81}
\newcommand{\LooseOverStructureN}{101}
\newcommand{\LooseOverStructureTrain}{85}
\newcommand{\LooseOverActivity}{66}
\newcommand{\LooseOverActivityN}{96}
\newcommand{\LooseOverActivityTrain}{65}
\newcommand{\LooseOverWorst}{81}
\newcommand{\LooseOverWorstN}{166}
\newcommand{\LooseOverWorstTrain}{77}
\newcommand{\LooseOverKing}{64}
\newcommand{\LooseOverKingN}{88}
\newcommand{\LooseOverKingTrain}{60}
\newcommand{\ActivityOverStructure}{75}
\newcommand{\ActivityOverStructureN}{24}
\newcommand{\ActivityOverStructureTrain}{82}
\newcommand{\ConflictTable}{%
\begin{table}[ht]\centering\small
\begin{tabular}{llrrrr}
\toprule
 & & \multicolumn{2}{c}{Train} & \multicolumn{2}{c}{Held out} \\
\cmidrule(lr){3-4}\cmidrule(lr){5-6}
First & against & $n$ & Sides with the first & $n$ & Sides with the first \\
\midrule
King safety & activity and space & 33 & 36 (22--53) & 32 & 44 (28--61) \\
King safety & pawn structure & 19 & 79 (57--91) & 9 & 56 (27--81) \\
King safety & something loose & 83 & 40 (30--51) & 88 & 36 (27--47) \\
King safety & the worst piece & 66 & 56 (44--67) & 55 & 64 (50--75) \\
Activity and space & pawn structure & 34 & 82 (66--92) & 24 & 75 (55--88) \\
Activity and space & something loose & 111 & 35 (27--44) & 96 & 34 (26--44) \\
Activity and space & the worst piece & 9 & 89 (56--98) & 8 & 75 (41--93) \\
Pawn structure & something loose & 114 & 15 (10--23) & 101 & 19 (12--28) \\
Pawn structure & the worst piece & 138 & 49 (41--58) & 120 & 50 (41--59) \\
Something loose & the worst piece & 163 & 77 (70--83) & 166 & 81 (75--87) \\
\bottomrule
\end{tabular}
\caption{Added after the run. The distinct positions with a clear verdict where both readings lean, in opposite directions: the share where the engine sides with the first, with 95\% Wilson intervals.}
\label{tab:conflict}
\end{table}
}
\newcommand{\MoveLogTrain}{86.1}
\newcommand{\MoveLogHeld}{87.7}
\newcommand{\MinLegal}{four}
\newcommand{\EaFen}{\fen{1r2k2r/p1p2ppp/4b3/2b1p3/3n1q2/3P1B1P/PPP2P2/RN1QK1NR w KQk - 0 13}}
\newcommand{\EaMover}{White}
\newcommand{\EaOther}{Black}
\newcommand{\EaPawns}{3.9}
\newcommand{\EaBetter}{Black}
\newcommand{\EaValue}{204}
\newcommand{\EaFavours}{Black}
\newcommand{\EaLine}{13.~b3 e4 14.~\symbishop{}xe4 \symknight{}xb3}
\newcommand{\EaBoard}{%
\chessboard[setfen={1r2k2r/p1p2ppp/4b3/2b1p3/3n1q2/3P1B1P/PPP2P2/RN1QK1NR w KQk - 0 13}, boardfontsize=14pt]
}
\newcommand{\EaWhiteMoves}{25}
\newcommand{\EaBlackMoves}{61}
\newcommand{\EaWhiteSpace}{5}
\newcommand{\EaBlackSpace}{20}
\newcommand{\EbFen}{\fen{8/8/4p1k1/p5pp/1P2N3/6K1/8/8 w - - 0 36}}
\newcommand{\EbMover}{White}
\newcommand{\EbOther}{Black}
\newcommand{\EbPawns}{5.9}
\newcommand{\EbBetter}{White}
\newcommand{\EbValue}{211}
\newcommand{\EbFavours}{Black}
\newcommand{\EbLine}{36.~bxa5 \symking{}f5 37.~a6 h4+}
\newcommand{\EbBoard}{%
\chessboard[setfen={8/8/4p1k1/p5pp/1P2N3/6K1/8/8 w - - 0 36}, boardfontsize=14pt]
}
\newcommand{\EbPassers}{three}
\newcommand{\EbMoverPassers}{zero}
\newcommand{\EbMoverIsolated}{one}


\title{Do the Judge the position factors tell who is better?\\[0.3em]
\large Five board readings against Stockfish, with material level, in 10-minute games at the
player's band}
\author{Chess Coach research lab}
\date{}

\begin{document}
\maketitle

\begin{abstract}
The app's Judge the position drill explains who is better with factors its board code scores in
centipawns. With material level, does each factor beyond material, and the worst piece, point to
the side Stockfish says is better? In \Unique{} positions from Lichess games at \Minutes{}
minutes between players at the player's band and the band above, all five do on both months, so
by the rule fixed before the run each \emph{tells who is better}. Ranked by held-out agreement:
\Ranked. All but pawn structure are right at least \ReliablePct\% of the time on their own.
Weighed together by a logistic model, the readings name the better side in \LogHeld\% of the
held-out positions with a clear verdict. Added after the run: whose move it is matters too.
Where nothing is loose, the side to move is the better side in \MoverHeldPct\%
of clear positions. A reading that favours the side to move is right \ForMoverMin\% to
\ForMoverMax\% of the time, while pawn structure and the worst piece, when they favour the side
that just moved, are no better than a coin. When two readings disagree, a loose pawn usually
wins and pawn structure usually loses.
\end{abstract}

\section{The question}

A coach judges a position before choosing a move: count the material, then look at the kings,
the pieces and the pawns. The app's Judge the position drill teaches that habit. It shows a
position from the player's game, asks who is better, and explains the engine's answer with five
factors the board code scores in centipawns: material, king safety, activity and space, pawn
structure, and something loose. The drill keeps a card only when the largest factor agrees with
the engine, so on its own cards a factor is right by construction. What its cards cannot show is
whether a factor points to the better side in the positions he will meet, which is what the habit
needs.

This study asks it of each factor beyond material, and of a sixth reading the app teaches
elsewhere, the worst piece, in positions where material is level. A reading that points to the
better side is worth teaching as a reason. One that does not would teach him to trust a sign that
tells him nothing, and one that points the wrong way is a bug in the board code.

\section{How it was measured}

\paragraph{The readings.} All are the app's, at commit \texttt{\AppCommit}, each from White's side
(positive when it favours White):

\begin{itemize}[nosep]
\item \textbf{King safety}: each king's danger is \WShield{} for each pawn missing from its
  shelter, \WRingAttacker{} for each point of weight of the enemy pieces attacking the squares
  around it, and \WOpenFile{} for each open file beside it (\WHalfOpenFile{} if half open); Black's
  danger minus White's, scaled down as the pieces come off.
\item \textbf{Activity and space}: \WMobility{} for each legal move more than the other side
  (counting each side's moves as if it were to move), and \WSpace{} for each square more it
  attacks in the other side's half.
\item \textbf{Pawn structure}: each side loses \WDoubled{} for a doubled pawn and \WIsolated{} for
  an isolated one, and gains \WPassed{} for a passed pawn and \WPassedRank{} for each rank it has
  advanced past its third, the passers counting up to twice as much as the pieces come off;
  White's sum minus Black's.
\item \textbf{Something loose}: the best capture balance of the side to move, less part of the
  other side's, each counted only when it wins material. In these positions, where neither side
  can win more than \TacticBar{} pawn by a capture, it leans only when the side to move can win a
  pawn and the other side has nothing to take, and it then favours the side to move.
\item \textbf{The worst piece}: the legal moves of White's least mobile piece (not a pawn or the
  king) minus those of Black's. ``Feed your laziest piece'' is one of the app's lessons; the drill
  folds it into activity.
\end{itemize}

A reading \emph{leans} when it is at least the drill's \LeanCp{} centipawns, or \WorstLean{}
moves for the worst piece, and it favours the side its sign names.

\paragraph{The games.} Lichess's rated games (the database's August 2026 file for training and
September 2026 held out) at \Minutes{} minutes with no increment, the player's time control, that
ended normally or on time after at least \MinPlies{} plies, with both players rated
\RatingLow{} to \RatingHigh: his band on Lichess (\BandLow) and the band above (\BandHigh). The
first \GamesPerMonth{} such games of each month enter.

\paragraph{The positions.} The drill's own population: after at least \PlyFloor{} plies, not over
and not in check, with at least \MinLegal{} legal moves, material exactly level, and no tactic
(no mate in one for either side, and no capture that wins more than \TacticBar{} pawn for either).
Either side may be to move. Of the \TrainLater{} positions after the first \PlyFloor{} plies in
the train games, \TrainQualifying{} (\TrainQualifyingPct\%) qualify, and \TrainLevelDrop\% fail because
material is not level. The held-out games give \HeldQualifying{} of \HeldLater{}
(\HeldQualifyingPct\%). For each reading, a game gives at most one position: one drawn at random,
seeded by the game and the reading, among its positions where the reading leans. Each reading
keeps up to \Cap{} positions a month, in file order: \Rows{} rows over the five readings,
\Unique{} distinct positions.

\paragraph{The verdict.} Stockfish through the app's engine service at the drill's grading budget
(\BudgetNodes{} nodes), with its hash cleared before each position and no tablebase. On the
drill's own scale, a position is \emph{level} within \LevelCp{} centipawns and \emph{clear} at
\ClearCp{} or more; a mate is clear. A reading's \emph{agreement} is the share of its positions
with a clear verdict where the side it favours is the side the engine says is better. Each
position comes from a different game, so its interval is the 95\% Wilson interval.

\paragraph{The rule, fixed before any position was scored.} For each reading, on each month
separately: it \emph{tells who is better} when its agreement's interval lies above 50\% on both;
it is \emph{reliable on its own} when the interval lies above \ReliablePct\% on both; it
\emph{points the wrong way} when the interval lies below 50\% on both; and there is \emph{no
signal} otherwise. Under \MinHeldout{} held-out positions with a clear verdict, the verdict is
\emph{not enough data}. Material is level in every position, so a reading that tells who is
better adds to material by construction.

\section{The verdict}

\RuleTable

All five readings tell who is better (Table~\ref{tab:rule}). On both months each agreement's
interval lies above 50\%, and the readings rank in the same order on both: \Ranked{} held out.
All but pawn structure are also reliable on their own, with intervals above \ReliablePct\%;
pawn structure falls short of that bar on both months. No reading points the wrong way, so the
board code has no factor whose sign is backwards.

\section{What the positions describe}

\DescTable

\paragraph{A lean is a sign, not a verdict.} The engine calls between \LevelMin\% and
\LevelMax\% of a reading's held-out positions level (Table~\ref{tab:desc}). The size of a lean
matters for the kings and the pieces: at twice its lean, king safety agrees in
\KingTwiceHeld\% of \KingTwiceHeldN{} held-out positions and activity in \ActivityTwiceHeld\% of
\ActivityTwiceHeldN, though these are few. Pawn structure and the worst piece gain nothing at
twice their lean. Something loose never reaches it here, since no capture wins more than a pawn.

\ResultTable

\paragraph{The result.} In the decisive games, the side a reading favours won \WonMin\% to
\WonMax\% of the time, on both months less often than the side the engine favours
(Table~\ref{tab:result}). Pawn structure trails the engine most: its side won \StructureWon\% of
these games, the engine's \StructureEngineWon\%.

\paragraph{Together.} The drill adds its factors as they come, in centipawns. That plain sum
favours the engine's side in \SumHeld\% of \SumHeldN{} held-out positions (\SumTrain\% on
train). A logistic model fitted on the train month, which weighs each reading, does a little
better: \LogTrain\% on train and \LogHeld\% held out. Per 100 centipawns it weighs something
loose most (\LooseWeight{} in log-odds), then activity (\ActivityWeight), king safety
(\KingWeight) and pawn structure (\StructureWeight). The worst piece, at \WorstWeight{} a move,
moves the model less at its lean than any drill factor does at its own, so beside the other four
it adds little.

\section{Whose move it is}

\SideTable

These descriptions were added once the verdicts were read; none was named before the run. They
change how the agreement above should be read.

\paragraph{The move.} Where nothing is loose, the side to move is the engine's better side in
\MoverHeld\% of the \MoverHeldN{} clear held-out positions (\MoverTrain\% on train). With material
level, a clear edge at this level is usually something the side to move can act on at once.
Every reading but something loose favours the side that just moved more often than the side to
move (Table~\ref{tab:side}), so for them the move works against agreement, not for it: the move
alone would put each below half. Something loose always favours the side to move, and agrees far
more often than the move alone would.

\paragraph{For the mover and against.} A reading that favours the side to move agrees
\ForMoverMin\% to \ForMoverMax\% of the time. One that favours the side that just moved still
agrees more often than not for king safety (\KingForOther\%) and activity (\ActivityForOther\%),
but pawn structure (\StructureForOther\%) and the worst piece (\WorstForOther\%) are then no
better than a coin on either month. When it is the only drill factor that leans, every reading
agrees well, pawn structure included (\StructureAlone\%).

\ConflictTable

\paragraph{When two disagree.} Where two readings lean opposite ways (Table~\ref{tab:conflict}),
the engine sides with something loose over activity (\LooseOverActivity\%), pawn structure
(\LooseOverStructure\%) and the worst piece (\LooseOverWorst\%) on both months, and over king
safety (\LooseOverKing\%) on the held-out month; and with activity over pawn structure
(\ActivityOverStructure\% of \ActivityOverStructureN). Pawn structure and the worst piece split
their disagreements. The other pairs are too few or too close to call.

\paragraph{The move as a reading.} With the side to move as a sixth input, the logistic model's
accuracy rises to \MoveLogTrain\% on train and \MoveLogHeld\% held out, and the move weighs more
than king safety at its lean and less than activity at its lean.

\section{Two positions}

Two held-out positions show a reading at work and the move outweighing one
(Figure~\ref{fig:examples}). They illustrate; they test nothing.

\begin{figure}[ht]
\centering
\begin{subfigure}[t]{0.48\textwidth}
\centering
\EaBoard\\[0.4em]
{\footnotesize FEN \EaFen}
\caption{\EaMover{} to move, in the middlegame. Only activity and space leans: \EaValue{}
centipawns for \EaFavours, who has \EaBlackMoves{} legal moves to \EaMover's \EaWhiteMoves{} and
attacks \EaBlackSpace{} squares in \EaMover's half to \EaWhiteSpace{} in \EaFavours's. Material is
level and no capture wins more than a pawn, yet the engine makes \EaBetter{} \EaPawns{} pawns better:
\EaLine.}
\label{fig:active}
\end{subfigure}\hfill
\begin{subfigure}[t]{0.48\textwidth}
\centering
\EbBoard\\[0.4em]
{\footnotesize FEN \EbFen}
\caption{\EbMover{} to move, in the endgame. Only pawn structure leans: \EbValue{} centipawns
for \EbFavours, who has \EbPassers{} passed pawns to none. But each side has a pawn to take and
\EbMover{} takes first, so the engine makes \EbBetter{} \EbPawns{} pawns better:
\EbLine.}
\label{fig:move}
\end{subfigure}
\caption{Held-out positions. Left: of the middlegame positions where activity and space is the
only drill factor that leans and the engine agrees with it, the one where it leans most. Right:
of the positions where pawn structure is the only drill factor that leans, favouring the side
that just moved, and the engine says the side to move is clearly better, the one where it leans
most. Each line is the engine's, at the grading budget.}
\label{fig:examples}
\end{figure}

\section{What the app can take}

\begin{itemize}[nosep]
\item \textbf{Keep the drill's factors.} Each points to the better side at his level, and none
  points the wrong way.
\item \textbf{Teach pawn structure as a tie-breaker.} On its own it is usually right, but when a
  loose pawn or the pieces say otherwise the engine sides with them, and when it favours the side
  that just moved it is no better than a coin.
\item \textbf{Name whose move it is.} The drill always shows a position with him to move, and the
  deck drops a card whose largest factor disagrees with the engine. Where the factors are quiet
  or favour his opponent, the engine often still says he is better because it is his move. The
  move as a sixth reason is a candidate for the drill, to be tested on his own cards with its
  rule fixed first.
\item \textbf{The worst piece is mostly activity.} Beside the four factors it adds little, so the
  drill is right to fold it in.
\end{itemize}

\section{Caveats}

\begin{itemize}[nosep]
\item \textbf{Not his games.} The positions are Lichess players' at his band and the band above;
  the drill's cards come from his own games.
\item \textbf{The engine's verdict, not his.} Whether a reading helps him choose moves is a later
  study's question, on his own answers.
\item \textbf{One reading at a time.} The readings move together, so a reading's agreement says it
  points the right way, not that it adds to the others. The logistic models describe that, and
  only describe.
\item \textbf{Added after the run.} The side to move, the readings alone and the disagreements
  were looked at once the verdicts were in. They describe these positions; a rule built on them
  needs positions this study did not see.
\item \textbf{Pawn breaks} are a plan rather than a standing edge, so they are tested as a rule of
  thumb, not here.
\end{itemize}

\end{document}
